% Noise Management Techniques in BGV-RNS?
% Different Variants of Key Switching?
% Gadget decomposition in BGV-RNS?
\section{Key-Switching}
\label{sec:gadget-decomposition}
\begin{itemize}
    \item A problem that arises in several areas of FHE is that of multiplying a ciphertext by some large number while keeping the error low.
    \item The techniques covered in this section typically occur in the context of keyswitching, but we remove them entirely from that context and consider only the techniques themselves.

    \item \textcolor{orange}{The gadget decomposition is not restricted to managing the noise in the scalar multiplication of ciphertexts, it is also central in the design of most FHE schemes as an auxiliary tool for certain FHE procedures; e.g., [BGV14, MP12, Bra12, GSW13, AP14, DM15, GINX16, CKKS17, CGGI20, BIP+22]} \cite{cryptoeprint:2012/078, cryptoeprint:2013/340, cryptoeprint:2014/816, CKKS, cryptoeprint:2018/421}
\end{itemize}

\begin{equation}
    (a\cdot s + \Delta m + \epsilon) \cdot X = aX\cdot s + \Delta Xm + X\epsilon
\end{equation}

% \textcolor{purple}{
% Several core operations require computing the product of a ciphertext element with noisy auxiliary material; performed naively, this scales the noise by a factor as large as $q$. This section covers the standard techniques for controlling this growth: decomposition (BV), modulus extension (GHS), and their combination (Hybrid).
% }

\newcommand{\decomp}[2]{\mathsf{Decomp}^{#2}(#1)}

\subsection{Brakerski-Vaikuntanathan}
Brakerski and Vaikuntanathan~\cite{cryptoeprint:2011/344} first proposed decomposing ciphertext elements into small digits as a means of controlling the noise growth \textcolor{red}{of this operation}. The idea has since been generalized into the idea of \emph{gadget decomposition} and applied in a variety of forms, including digit and CRT (RNS) decomposition. Following the convention of~\cite{cryptoeprint:2021/204, cryptoeprint:2022/915}, we refer to this family of decomposition-based techniques collectively as \emph{BV}.

\begin{definition}[Digit Decomposition]
    Given an integer $\gamma\in\mathbb{Z}_q$ for a modulus $q\geq 2$, base $\beta$ and decomposition parameter $\ell$, the digit decomposition $\decomp{\gamma}{} = (\gamma_0, \gamma_1, \dots, \gamma_{\ell-1}) \in \mathbb{Z}_\beta^\ell$ outputs an ordered list of $\ell$ base-$\beta$ digits satisfying
    \begin{equation*}
        \gamma = \gamma_0\frac{q}{\beta} + \gamma_1\frac{q}{\beta^2} + ... \gamma_{\ell - 1}\frac{q}{\beta^{\ell}} = \sum_{i=0}^{\ell-1}\gamma_i\frac{q}{\beta^{i+1}}
    \end{equation*} 
\end{definition}

%%
%% Why its helpful
Digit decomposition has two important properties. First, if the digits $\gamma_i$ of $\decomp{\gamma}{}$ are small, then $\langle \decomp{\gamma}{}, \epsilon\rangle \leq \ell\cdot \max\{\gamma_i\cdot \epsilon_i\} \leq \beta \epsilon \ll \gamma \epsilon$. This is the property that allows some control over the noise growth. Second, two integers can be multiplied through the decomposition. By construction, $\langle\decomp{\gamma}{}, g\rangle = \gamma$ for the fixed vector $g = (q/\beta, \dots, q/\beta^{\ell})$, meaning $\langle\decomp{\gamma}{}, g\cdot x\rangle = \gamma\cdot x$ for an arbitrary integer $x\in \mathbb{Z}$. For either of these properties, the particular form of $g$ is unimportant as long as the digits $\gamma_i$ remain small. Furthermore, as we are generally working with some intrinsic error in FHE, \cite{cryptoeprint:2018/421} proposed allowing some error $\epsilon$ in the reconstruction, giving the definition from their paper verbatim:

%%
%% Gadget Decomposition
\begin{definition}[(Abstract) Gadget Decomposition]
    \label{def:tfhe:gadget-decomposition}
    An efficient algorithm $\decomp{v}{g,\beta,\epsilon}$ is a valid decomposition on the gadget $g\in \mathbb{Z}_q^\ell$ with quality $\beta \in \mathbb{Z}^{+}$ and precision $\epsilon \in \mathbb{Z}^+$ if and only if for any \gls{rlwe} sample $v\in R_Q$ it efficiently and publicly outputs a small vector $u \in \mathbb{Z}_q$ such that $\|u\|_\infty \leq \beta$ and $\langle u, g\rangle \leq v + \epsilon$.

    % Furthermore, the expectation of u · H − v must to be equal to 0 when v is uniformly distributed in M .
\end{definition}

\begin{itemize}
    \item Other (equivalent) form more convenient to work with, where the output of $\decomp{v}{g,\beta,\epsilon} \to D(v)$ and $g\cdot u \to P(u)$ are considered.
\end{itemize}

\begin{definition}[Gadget Property]
    Two vectors $P(x) \in \mathbb{Z}_Q^\ell$ and $D(y) \in \mathbb{Z}_\beta^\ell$ are said to have the \textit{gadget property} if and only if they satisfy \eqref{def:gadget-property} for any pair of \gls{rlwe} samples $(x,y)\in R_Q^2$ with precision $\epsilon \in \mathbb{Z}_q^+$ and quality $\beta$ where $\|D(y)\|_\infty \leq \beta$. If $\epsilon = 0$ then $P(x)$ and $D(y)$ are said to have the \textit{exact} gadget property.
    \begin{equation}
        \label{def:gadget-property}
        \langle P(x), D(y)\rangle = x\cdot y + \epsilon \pmod{Q}
    \end{equation}
\end{definition}

% \begin{itemize}
%     \item Can be applied component-wise to polynomials $\gamma\in R_q$ instead of $\mathbb{Z}_q$
%     \item Centered representation
% \end{itemize}

\begin{lemma}%[Layered Gadget Decomposition]
    \label{thm:gadget-layers}
    The composition of two gadget decompositions is itself a valid gadget decomposition.
    \begin{gather*}
        \langle P_1(u), D_1(v)\rangle = u\cdot v + \epsilon_1, \quad \langle P_2(u), D_2(v) \rangle = u\cdot v + \epsilon_2 \pmod{Q} \\
        P_1\circ P_2 = P_1(P_2(u)) \implies \langle P_1\circ P_2, D_1 \circ D_2 \rangle = u\cdot v + 2\epsilon \pmod{Q} 
    \end{gather*}

    \begin{proof}
        This is clear from applying the gadget property twice.
        \begin{gather*}
            \langle P_1(P_2(u)), D_1(D_2(u)) \rangle = P_2(u) \cdot D_2(u) + \epsilon_1 = u\cdot v + \epsilon_2 + \epsilon_1 \pmod{Q}
        \end{gather*}
    \end{proof}
\end{lemma}






\subsubsection{RNS Instantiation}
Digit decomposition cannot be applied directly the \gls{rlwe} sample in an \gls{rns} scheme, where the system only has access to the residues. When Bajard et. al. first implemented the BV cryptosystem using \gls{rns} \cite{cryptoeprint:2016/510} they exchanged digit decomposition for \gls{crt} decomposition, exploiting the fact that this decomposition is already required by the \gls{rns}.

%% TODO: Clean up, make sure \hat{q} notation is known 
By stating the \gls{crt} as an inner product \autoref{eq:crt-gadget}, it is immediately clear that \gls{crt} with moduli $q_i$ is an exact decomposition on $g = (\hat{q}_0\cdot\hat{q}_0^{-1}\bmod{Q}, \dots, \hat{q}_{k-1}\cdot\hat{q}_{k-1}^{-1}\bmod{Q})$ with $\beta = \max_i\|q_i\|_\infty/2$ (assuming centered residues) and $\epsilon = 0$. 

\begin{itemize}
    \item CRT idempotent used in \cite{cryptoeprint:2016/510}
    \item HPS optimization for specifically keyswitching \cite{cryptoeprint:2018/117}
    \item Definition: gadget with $\beta = Q/2, \epsilon = 0$
\end{itemize}

\begin{align}
    \hps{P}(u) = () \\
    \hps{D}(v) = ()
\end{align}

\begin{itemize}
    \item Note on HPS gadget = RNS limb lookup $\to$ \autoref{sec:impl:core}
    \item HPS is exact but noise is large ($\beta = Q/2,\; \epsilon = 0$); can apply digit decomposition to each residue polynomial as per \autoref{thm:gadget-layers}, lowering $\beta \to B/2$ and still exact ($\epsilon = 0$) if $B = \lfloor \max{q_i}/\ell \rfloor + 1$
\end{itemize}
% \subsection{Gentry-Halevi-Shoup}
\begin{itemize}
    \item Another approach was implemented by Gentry, Halevi and Shoup in \cite{cryptoeprint:2012/099}. 
    \item Generate the key in $QP$ as $(a\cdot s + \Delta + \epsilon, a) \in \mathcal{B}_{QP}^2$
    \item Base extend $u \mapsto \mathcal{B}_{QP}$ and calculate inner product $\rlwe{(mP)}$
    \item Divide the ciphertext by $P$ and approx mod down to $\mathcal{B}_{Q}$
    \item \textcolor{red}{\textbf{FastBaseExtension and Multiply by P = add empty towers!!!}}
    \item Since key was pre-scaled, $\epsilon$ shrinks relative to the message
    \item Reference RNS section, base conversion
\end{itemize}

\begin{remark}[GHS Security] % In RGSW, $Z$ provides the security of the scheme by masking $m$. 
    \label{remark:ghs-security}
    The public key is masked by elements in $QP$, meaning the ciphertext modulus $QP \approx 2Q$ must be used when evaluating the security. A GHS scheme has the security as a BV scheme if $Q$ is halved. 
\end{remark}

\begin{lemma}
    GHS is a single-digit decomposition $(\ell = 1)$ with quality $\beta = Q/P$ and precision $\epsilon = 0$. 
    \begin{gather} %%%% Exact gadget
        \langle [x]_{\mathcal{B}_{QP}}, P\cdot [y]_{\mathcal{B}_{QP}}\rangle = x\cdot y \cdot P \pmod{QP} \\
        \langle P^{-1}\cdot [x]_{\mathcal{B}_{QP}}, P\cdot [y]_{\mathcal{B}_{QP}}\rangle = x\cdot y \cdot P \cdot P^{-1} = x\cdot y \pmod{QP} \\
        \overset{\tiny\text{ModDown}}{\longmapsto} \;\; x\cdot y \pmod{Q}
    \end{gather}

    %%%% Inexact gadget
    % \begin{gather}
    %     \langle [x]_{\mathcal{B}_{QP}}, P\cdot [y]_{\mathcal{B}_{QP}}\rangle = x\cdot y \cdot P + \epsilon \pmod{QP} \\
    %     \implies \langle P^{-1}\cdot [x]_{\mathcal{B}_{QP}}, P\cdot [y]_{\mathcal{B}_{QP}}\rangle = x\cdot y \cdot P \cdot P^{-1} + \epsilon = x\cdot y + \epsilon \pmod{QP} \\
    %     \overset{\tiny\text{ModDown}}{\longmapsto} x\cdot y + \epsilon/P \pmod{Q}
    % \end{gather}
\end{lemma}

\subsection{Hybrid}
\label{sec:hybrid}
\begin{itemize}
    \item BV and GHS represent a spectrum
    \item Hybrid = combination of the two techniques
    \item Group digits by $d_{num}$
\end{itemize}
\subsection{RGSW}
The techniques above are used in the \gls{gsw} scheme~\cite{cryptoeprint:2013/340} and its ring variants~\cite{cryptoeprint:2014/816, cryptoeprint:2018/421}, where the ciphertext itself is built around a gadget. We now formally define an \gls{rgsw} ciphertext, and the operations required to employ \gls{rgsw} in \gls{fhe} schemes. %% TODO: Rewrite for style + TFHE approximates gadget + Formalize sample vs ciphertext vs encryption.

\begin{definition}[Gadget Matrix]
    \label{def:gadget-matrix}
    A gadget $g \in R_Q^\ell$ is extended to $k$-dimensional ciphertexts by the block-diagonal gadget matrix $G = I_k \otimes g \in R_Q^{k\ell\times k}$
\end{definition}

\begin{definition}[RGSW Ciphertext]
    Given a gadget matrix $G = I_2 \otimes g$, we define the \emph{masking} matrix $Z \in R_Q^{2\ell\times 2}$ as a collection of $2\ell$ \gls{rlwe} public keys $\Call{RLWE}{0}$. An \emph{\gls{rgsw}} encryption $C$ of $\mu \in R_Q$ is denoted $\Call{RGSW}{\mu}$ and defined as
    \begin{equation}
        \label{def:rgsw}
        C = Z + \mu\cdot G % \in R_Q^{2\ell\times 2}
    \end{equation}
\end{definition}

% Each row of an \gls{rgsw} ciphertext is a valid \gls{rlwe} sample. The upper block of $\ell$ rows are equivalent to $\Call{RLWE}{\mu g_i}$ for row $i$. The lower block of the remaining $\ell$ rows encode the message in the public components, where the $i$-th row effectively becoming $\Call{RLWE}{-\mu g_i\cdot s}$.

An \gls{rgsw} ciphertext encodes the gadget-scaled messages $\mu g_i$ row by row. In the top block, $\mu g_i$ is added to the masked element, replacing the zero message of the underlying $\Call{RLWE}{0}$; in the bottom block it is added to the public element, so the row decrypts to $-\mu g_i\cdot s$ instead. Every row of $C$ is therefore itself an \gls{rlwe} ciphertext.

Two \gls{rgsw} ciphertexts add homomorphically, just as \gls{rlwe} ciphertexts do: the zero encryptions and the gadget terms add component-wise, so $C_A + C_B$ encrypts $\mu_A + \mu_B$. Addition alone is, however, not sufficient for \gls{fhe}. This motivates the need for an \gls{rgsw}$\times$\gls{rgsw} product and a \gls{rlwe}$\times$\gls{rgsw} product. These are called the internal and external products, respectively, in \cite{cryptoeprint:2018/421}. First, we extend the gadget property to two dimensions:
% The multiplication mirrors the vector--matrix relationship between the two ciphertext types: the \gls{rlwe} ciphertext $x = (x_0, x_1)$ is mapped by a two-dimensional decomposition $G^{-1}(x)$ to a row vector of $2\ell$ small ring elements, which multiplies the $2\ell\times 2$ matrix $C$. The only requirement on $G^{-1}(x)$ is that it approximately inverts the $G$ matrix: 
\begin{equation}
    \label{eq:gadget-matrix-property}
    G^{-1}(x)\cdot G = x + e \pmod{Q}%, \qquad\|G^{-1}(x)\|_{\infty} \leq \beta
\end{equation}

%% TODO: Rewrite if time
No particular decomposition is prescribed. The canonical choice is the base-$B$ digit decomposition of the ciphertext coefficients~\cite{cryptoeprint:2018/421}, and it is the one usually meant when the external product is written down, but nothing in the definition depends on it. Any $G^{-1}(x)$ meeting \eqref{eq:gadget-matrix-property} at quality $\beta$ yields a valid product, and the noise bounds below are stated in terms of $\beta$ alone.

\begin{definition}[External Product]
    \label{def:prod:external}
    Let $x \in R_Q^2$ be an \gls{rlwe} encryption of $m \in R_t$, let $C \in R_Q^{2\ell\times 2}$ be an \gls{rgsw} encryption of $\mu \in R_Q$ under the same secret key, and let $G^{-1}$ be a gadget decomposition satisfying \eqref{eq:gadget-matrix-property}. The {external product} is defined as
    \begin{align}
        \boxdot: \Call{RLWE}{} \times \Call{RGSW}{} &\longmapsto \Call{RLWE}{} \nonumber \\
        \label{def:extprod}
        (x, C) &\longmapsto x \boxdot C = G^{-1}(x) \cdot C,
    \end{align}
    which evaluates $\Call{RLWE}{m}\,\boxdot\,\Call{RGSW}{\mu} = \Call{RLWE}{m\cdot\mu}$. Formal correctness and noise bounds are given in~\cite{cryptoeprint:2018/421}, though informally this follows from inserting \eqref{def:rgsw} into \eqref{def:extprod} and using the gadget property \eqref{eq:gadget-matrix-property}, yielding a sum of two \gls{rlwe} samples:
    \begin{equation}
        \label{eq:extprod-expansion}
        G^{-1}(x)\cdot C = G^{-1}(x)\cdot Z + \mu\cdot(x + e)
    \end{equation}
\end{definition}

\vspace{1pt}

\begin{lemma}[Gadget Matrix Construction]
    \label{lemma:gadget-matrix-construction}
    Two vectors $P(u), D(v)$ that satisfy the gadget property \eqref{def:gadget-property} define a valid pair of gadget matrices \eqref{def:Gmat} and \eqref{def:Ginvmat} satisfying \eqref{eq:gadget-matrix-property}, with $\|G^{-1}(x)\|_\infty \leq \beta$ inherited from the quality of $D(v)$.
    \begin{gather}
        \label{def:Gmat}
        G = I_2 \otimes P(1) \\
        \label{def:Ginvmat}
        G^{-1}(x) = [D(x_0)|D(x_1)]
    \end{gather}

    \begin{proof}
        The block structure reduces the product to one inner product per component, each an instance of \eqref{def:gadget-property}:
        \begin{gather*}
            G^{-1}(x)\cdot G = [D(x_0)|D(x_1)]
            \begin{pmatrix}
                P(1) & 0 \\
                0 & P(1)
            \end{pmatrix}, \quad \|D(\cdot)\|_\infty \leq \beta \\
            \big(\langle D(x_0), P(1)\rangle,\; \langle D(x_1), P(1)\rangle\big)
            = (x_0 + e_0,\; x_1 + e_1) = x + e \pmod{Q}
        \end{gather*}
    \end{proof}
\end{lemma}

\begin{remark}%[Asymmetric Noise Growth]
    \label{rem:extprod-asymmetry} %% TODO: Include more sources/authors
    \gls{rgsw} ciphertexts have been by several authors, including \cite{cryptoeprint:2013/541, cryptoeprint:2018/421}, to be uniquely suited to evaluate long chains of products, due to an asymmetry in the noise propagation --- an asymmetry explained in greater detail in \cref{appendix:asymmetry}. Essentially, since $\|G^{-1}(x)\|_\infty \leq \beta$ and $\mu\in\{0,1\}$ we see that \eqref{eq:extprod-expansion} is bounded by $\beta\cdot Z + x + e$. The total error in the external product becomes dominated by the product $\beta \cdot Z$, while $x$ is effectively ignored. A chain of external products $(x\boxdot C_0) \boxdot C_1 \boxdot ...$ is thus able to sustain itself with minimal noise growth as long as the \gls{rgsw} ciphertexts $C_i$ are fresh encryptions.
\end{remark}

\begin{definition}[Internal Product]
    \label{def:prod:internal}
    Let $A, C \in R_Q^{2\ell\times 2}$ be \gls{rgsw} encryptions of $\mu_A, \mu_C \in R_Q$ under the same secret key and gadget, and let $A_i$ denote the $i$-th row of $A$. The \emph{internal product} applies the external product to each row of the left operand:
    \begin{align}
        \boxtimes: \Call{RGSW}{} \times \Call{RGSW}{} &\longrightarrow \Call{RGSW}{} \nonumber \\
        (A, C) &\longmapsto A \boxtimes C = \begin{bmatrix}
            A_0 \boxdot C \\
            A_1 \boxdot C \\
            \vdots \\
            A_{2\ell - 1} \boxdot C
        \end{bmatrix},
    \end{align}
    which returns an encryption of the product of the messages: $\Call{RGSW}{\mu_A} \boxtimes \Call{RGSW}{\mu_C} = \Call{RGSW}{\mu_A\cdot\mu_C}$.
\end{definition}

One internal product amounts to $2\ell$ external products, and correctness follows row by row: each external product multiplies the content of $A_i$ by $\mu_C$, so the stacked result is again of the form \eqref{def:rgsw} with message $\mu_A\cdot\mu_C$.

The asymmetry of \autoref{rem:extprod-asymmetry} applies to each of the $2\ell$ rows, and dictates the order in which a chain of products should be evaluated. The noise of the right-hand operand is amplified by the decomposition in every row of the result, while the noise of the left-hand operand is only carried through. When accumulating a running product, the accumulator must therefore be kept on the left with each fresh ciphertext supplied on the right, so that the noise grows additively rather than multiplicatively.

\begin{remark}[Operand Order]
    \label{rem:operand-order}
    The convention above places the \gls{rlwe} operand of the external product on the left and the \gls{rgsw} operand on the right, so that both products read left to right with the accumulator first. TFHE writes the external product the other way round~\cite{cryptoeprint:2018/421}, with the \gls{rgsw} ciphertext on the left. The operation is the same, but any expression carried over from that literature must be flipped before it is read against \eqref{def:extprod}.
\end{remark}
